\documentclass[10pt, twocolumn]{article}

\usepackage[margin=1.8cm]{geometry}
\usepackage{amsmath, amssymb}
\usepackage{graphicx}
\usepackage{booktabs}
\usepackage{hyperref}
\usepackage[numbers]{natbib}
\usepackage{microtype}
\usepackage{caption}
\usepackage{subcaption}
\usepackage{float}
\usepackage{placeins}
\usepackage{lmodern}
\usepackage[T1]{fontenc}
\usepackage[utf8]{inputenc}
\usepackage{balance}

\captionsetup{font=small, labelfont=bf}
\hypersetup{colorlinks=true, linkcolor=blue, citecolor=blue, urlcolor=blue}

% -----------------------------------------------------------------------
\title{Predicting the Rating of Perceived Exertion in Competitive
       Runners Using Tree-Based Machine Learning Models}

\author{Lucas Savona\\
        \small Master IEAP DigiMove\\
        \small Universit\'{e} de Montpellier\\\\
        \small \url{https://github.com/lucassavona85-png/Tree_based_RPE_Prediction}}

\date{April 2026}
% -----------------------------------------------------------------------

\begin{document}

\maketitle

% -----------------------------------------------------------------------
\begin{abstract}
Accurately predicting an athlete's subjective training load the day after
a session can inform training management and reduce injury risk. This study
investigates whether the Rating of Perceived Exertion (RPE) can be predicted
from the preceding six days of training history using tree-based machine
learning models. Three models, a Decision Tree, a Random Forest, and
XGBoost, are trained and evaluated on a longitudinal dataset of 74
competitive runners \citep{lovdal2021}. Two validation strategies are
compared: a single chronological 80/20 split and a five-fold walk-forward
cross-validation scheme. XGBoost achieves the best performance under both
schemes (RMSE\,=\,0.239, MAE\,=\,0.189, $R^2$\,=\,0.239 on the held-out
test set), with past training success emerging as the most informative
predictor. Results confirm the viability of tree-based approaches for
personalised training load monitoring.
\end{abstract}

% -----------------------------------------------------------------------
\section{Introduction}
\label{sec:intro}

The goal of this study is to predict the Rating of Perceived
Exertion (RPE) of a competitive runner's training session from their
recent training history covering the preceding six days. RPE is a
subjective measure of exercise intensity that integrates physiological
and psychological load into a single scalar value \citep{borg1982, foster2001}.
Reliable RPE forecasts could allow practitioners to anticipate excessive
training stress and intervene before fatigue accumulates into injury.

Predicting RPE from training data has been explored in team sports
contexts. L{\"o}vdal et al.~\cite{lovdal2021} demonstrated that XGBoost models trained
on GPS and subjective load data from competitive runners could predict
injury risk with an AUC of 0.724. Related work in football has shown
that external load variables measured over rolling windows can predict
session RPE with moderate accuracy, though results depend heavily on
the athlete population and the prediction horizon \citep{vallance2023}.
To our knowledge, no study has framed RPE prediction as a pure
time-series regression task on this dataset, using only the preceding
six days of history as input while explicitly excluding same-session
features to avoid data leakage.

Machine learning offers a principled framework for this prediction task.
Tree-based models are particularly well suited because they (i)~handle
multivariate, heterogeneous feature sets without preprocessing, (ii)~are
inherently interpretable through variable importance measures, and (iii)~have
demonstrated competitive performance in sports science applications
\citep{lovdal2021}. We evaluate three such models of increasing complexity:
a single Decision Tree (CART), a Random Forest ensemble, and a gradient
boosting model (XGBoost).

Each model is selected to represent a distinct point on the
interpretability-performance spectrum. The Decision Tree offers full
transparency: its structure can be visualised as a single diagram
(Figure~\ref{fig:tree}) and every prediction traced to a root-to-leaf
path. The Random Forest sacrifices individual tree transparency for
lower variance through bagging. XGBoost further reduces bias through
sequential boosting at the cost of additional complexity. This
progression allows us to empirically quantify the accuracy gain
achieved at each step of the ensemble hierarchy.

A critical methodological challenge in athlete monitoring is temporal
dependency: training load and recovery follow temporal patterns, so a naive
cross-validation that shuffles observations would introduce data leakage
across time. We address this by using chronological validation throughout
and compare two schemes, a single temporal split and a walk-forward
cross-validation, to assess their consistency.

% -----------------------------------------------------------------------
\section{Tree-Based Models}
\label{sec:models}

\subsection{Decision Trees (CART)}
\label{subsec:dt}

A Classification and Regression Tree (CART) \citep{breiman1984} partitions
the feature space by recursively selecting the split variable and threshold
that maximally reduce the mean squared error (MSE) of the response within
each resulting region. Formally, at node~$t$ the algorithm selects the
split that maximises the impurity decrease:

\begin{equation}
  \Delta_{\text{MSE}}(t) = \text{MSE}(t)
    - \frac{|t_L|}{|t|}\,\text{MSE}(t_L)
    - \frac{|t_R|}{|t|}\,\text{MSE}(t_R),
  \label{eq:mse_split}
\end{equation}

where $t_L$ and $t_R$ denote the left and right child nodes and
$|\cdot|$ their cardinality. The tree grows until a stopping criterion
is met, controlled by the complexity parameter $cp$ (minimum impurity
decrease required for any split). In this study, $cp$ is treated as
the primary hyperparameter and selected by chronological cross-validation.

Unlike linear models, decision trees require neither feature normalisation
nor dimensionality reduction, as the splitting criterion implicitly selects
the most informative variables at each node. The final tree is visualised
using a recursive node diagram (Figure~\ref{fig:tree}), providing a fully
transparent decision path from raw features to the predicted RPE value.

A known limitation is that predictions in leaves populated by few training
examples carry higher uncertainty, as the leaf mean is estimated from a
small sample \citep{breiman1984}. This motivates the ensemble methods
described below. In a related line of research, Sutton-Charani~\citep{suttoncharani2014}
extended CART to uncertain training data by modelling imprecise attributes
and labels with belief functions and estimating tree parameters via the
credibilistic E$^2$M algorithm.

\subsection{Random Forests}
\label{subsec:rf}

A Random Forest \citep{breiman2001} mitigates the high variance of individual
trees by averaging the predictions of $B$ independently trained trees:

\begin{equation}
  \hat{y} = \frac{1}{B} \sum_{b=1}^{B} T_b(\mathbf{x}),
  \label{eq:rf_avg}
\end{equation}

where each tree $T_b$ is trained on a bootstrap sample of the training data
and, at each split, considers only a random subset of $m$ features
(we use $m = \lfloor p/3 \rfloor$ for regression, following the standard
default \citep{hastie2009}). This \emph{decorrelation} between trees reduces
variance without substantially increasing bias, yielding lower generalisation
error than a single tree.

Variable importance is estimated via permutation: for each feature, the
increase in out-of-bag prediction error when its values are randomly
shuffled is recorded (\%IncMSE). This permutation-based measure is
preferred over the mean decrease in node impurity, which is biased towards
high-cardinality features \citep{breiman2001}.

The bias-variance decomposition underpins the expected performance
progression: a single deep tree has low bias but high variance;
Random Forest reduces variance through bagging; gradient boosting
reduces bias through sequential residual fitting while controlling
variance via shrinkage ($\eta$) and sub-sampling.

\subsection{Gradient Boosting (XGBoost)}
\label{subsec:xgb}

Gradient boosting \citep{friedman2001} builds an additive ensemble by
sequentially fitting shallow trees to the negative gradient of the loss
function:

\begin{equation}
  F_m(\mathbf{x}) = F_{m-1}(\mathbf{x}) + \eta \cdot h_m(\mathbf{x}),
  \label{eq:gb_update}
\end{equation}

where $h_m$ is a regression tree fitted to the residuals of $F_{m-1}$
and $\eta \in (0,1]$ is the learning rate (also called \emph{shrinkage}).
The learning rate is the most critical hyperparameter: a smaller $\eta$
requires more boosting rounds but typically achieves better generalisation.
XGBoost \citep{chen2016} extends this framework with second-order gradient
approximations, column sub-sampling, and $\ell_1$/$\ell_2$ regularisation,
making it both faster and more robust than vanilla gradient boosting.
Feature importance is measured by the \emph{Gain} criterion: the average
reduction in training loss attributable to each feature across all splits
where it appears.

% -----------------------------------------------------------------------
\section{Data and Methods}
\label{sec:methods}

\subsection{Dataset}
\label{subsec:data}

We use the \emph{Injury Prediction for Competitive Runners} dataset
collected by L{\"o}vdal et al.~\cite{lovdal2021}, publicly available on
Kaggle\footnote{\url{https://www.kaggle.com/code/mohamedbakrey/anticipate-injured-numbers-in-competitive-runners}}.
It comprises seven years of daily training logs from 74 high-level middle- and
long-distance runners. For each athlete and day, ten variables are recorded:
total kilometres run, kilometres in heart-rate zones 3--4 and 5/T1/T2,
sprint kilometres, number of sessions, number of strength sessions, RPE,
perceived recovery, and a binary injury indicator.

Our target variable is RPE, normalised to the interval $[0, 1]$. Following
L{\"o}vdal et al.~\cite{lovdal2021}, rows where RPE is missing are removed; the target is
never imputed. The original binary injury label is not used. After filtering,
the modelling dataset contains 42\,766 observations.

Table~\ref{tab:features} summarises the ten raw features recorded
daily for each athlete. The target variable (RPE, no suffix) exhibits
a right-skewed distribution with mean 0.22 and standard deviation 0.14
on the training set, reflecting the predominance of low-intensity
sessions in competitive runners' training programmes.

\begin{table}[!htb]
  \centering
  \small
  \caption{Raw daily features in the dataset. Each feature is
           replicated six times (suffixes .1--.6) to form the
           60-column feature matrix used for modelling.}
  \label{tab:features}
  \begin{tabular}{@{}llp{3.8cm}@{}}
    \toprule
    Feature & Type & Description \\
    \midrule
    Total km           & Obj. & Total distance run \\
    km Z3-4            & Obj. & Distance in HR zones 3--4 \\
    km Z5-T1-T2        & Obj. & Distance in HR zones 5/T1/T2 \\
    km sprinting       & Obj. & Sprint distance \\
    Nr.\ sessions      & Obj. & Sessions per day \\
    Strength training  & Obj. & Strength sessions \\
    Hours alternative  & Obj. & Cross-training hours \\
    Perceived exertion & Subj. & RPE (target at day $t$) \\
    Training success   & Subj. & Self-rated session quality \\
    Perceived recovery & Subj. & Self-rated recovery \\
    \bottomrule
  \end{tabular}
\end{table}

\subsection{Feature Engineering}
\label{subsec:features}

The dataset is provided as a pre-built sliding window: each row already
contains the current day's variables (suffix-free, day~$t$) alongside the
values of all ten variables for the six preceding days (suffixes \texttt{.1}
through \texttt{.6}, corresponding to days $t-1$ through $t-6$). An
empirical shift test confirmed that the no-suffix RPE column corresponds
to the current day and must serve as the target, never as a feature.

The 60 feature columns (10 variables $\times$ 6 lag days) are used directly
as model inputs. No normalisation or dimensionality reduction is applied;
as noted in Section~\ref{sec:models}, tree-based models are invariant to
monotone transformations of features.

\subsection{Experimental Protocol}
\label{subsec:protocol}

Train/test split.
A chronological split is used: observations with a global date index at or
below the 80th percentile (Date~$\leq 2057$) constitute the training set
(34\,226 observations) and the remaining 20\% form the held-out test set
(8\,540 observations). To preserve the temporal integrity of the data,
no shuffling is applied, as randomised cross-validation would introduce
data leakage across athletes over time.

Validation schemes.
Two schemes are compared, applied to the training set only; the held-out
test set is never seen during training or hyperparameter selection.

\begin{itemize}
  \item Temporal Split (TS): Models are fitted once on the full
    training set after an internal time-slice cross-validation selects the
    optimal hyperparameters. Performance is reported once on the test set.
  \item Walk-Forward CV (FW): The training set is divided into five
    non-overlapping folds using an expanding-window strategy
    (initialWindow\,=\,70\%, horizon\,=\,6\% of training rows,
    fixedWindow\,=\,FALSE). Each fold evaluates predictions on a block of
    future observations not seen during training. The mean $\pm$ standard
    deviation of the fold-level RMSE is reported for model selection.
    The winning model is then re-fitted on the full training set and
    evaluated on the shared test set.
\end{itemize}

Hyperparameters.
For the Decision Tree, the complexity parameter $cp \in [0.0001, 0.05]$
is selected by minimising RMSE over the walk-forward folds (or the
internal time-slice CV for the TS scheme); the optimal value is
$cp = 0.001821$, yielding a tree with 17 leaves. For the Random Forest,
500 trees are grown with $m = 20$ candidates per split. For XGBoost,
$\eta = 0.05$, max depth~=~4, column sub-sampling~=~0.8, and the number
of boosting rounds (527) is determined by five-fold early stopping
(50 rounds patience).

Baseline.
A dummy predictor that constantly outputs the training-set mean RPE (0.223)
is included as a performance floor. Any model with $R^2 \leq 0$ fails to
outperform this trivial baseline.

Evaluation metrics.
Three complementary metrics are reported on the held-out test set:
Root Mean Squared Error (RMSE, penalises large errors), Mean Absolute Error
(MAE, robust to outliers), and the coefficient of determination $R^2$
(proportion of variance explained, model-agnostic).

\subsection{Descriptive Analysis of the Target Variable}
\label{subsec:eda}

\begin{figure}[H]
  \centering
  \includegraphics[width=0.95\columnwidth]{../figures/fig_rpe_histogram.pdf}
  \caption{Distribution of RPE on the training set ($n = 34\,226$).
           Dashed red: mean; dotted blue: median. The distribution
           is right-skewed with mass between 0.1 and 0.3.}
  \label{fig:rpe_hist}
\end{figure}

Figure~\ref{fig:rpe_hist} shows that RPE is right-skewed, with mass
concentrated between 0.1 and 0.3, reflecting the predominance of
low-intensity sessions; high-intensity sessions (RPE $> 0.5$) are
rare. The dummy baseline consequently over-predicts low-intensity
sessions and under-predicts high-intensity ones ($R^2 = -0.206$).

\begin{figure}[H]
  \centering
  \includegraphics[width=0.95\columnwidth]{../figures/fig_correlation_heatmap.pdf}
  \caption{Pairwise Pearson correlation at lag $t-1$.
           Subjective variables cluster together; objective and
           subjective features form distinct correlation blocks.}
  \label{fig:corr_heatmap}
\end{figure}

Figure~\ref{fig:corr_heatmap} reveals that subjective variables
(exertion, training success, recovery) are moderately inter-correlated
but weakly correlated with objective load metrics, suggesting
complementary dimensions of the training response. Because athletes
vary in absolute volumes, these implicitly self-normalised subjective
indicators may reduce inter-individual variability, a hypothesis
supported by the feature importance rankings
(Section~\ref{sec:results}).

\FloatBarrier
% -----------------------------------------------------------------------
\section{Results}
\label{sec:results}

\subsection{Walk-Forward Cross-Validation}

Table~\ref{tab:cv} reports the fold-level RMSE for the five walk-forward
folds under the FW scheme. XGBoost achieves the lowest mean CV RMSE, with
small standard deviation, indicating stable performance across time windows.
The CV ranking (DT $>$ RF $>$ XGBoost) is consistent across all five folds
(Figure~\ref{fig:cv_folds}, Appendix).

\begin{table}[htbp]
  \centering
  \small
  \caption{Walk-forward cross-validation RMSE (mean\,$\pm$\,sd, 5 folds).}
  \label{tab:cv}
  \begin{tabular}{lc}
    \toprule
    Model & CV RMSE \\
    \midrule
    Decision Tree & $0.2534 \pm 0.0105$ \\
    Random Forest & $0.2444 \pm 0.0084$ \\
    XGBoost       & $\mathbf{0.2416 \pm 0.0080}$ \\
    \bottomrule
  \end{tabular}
\end{table}

\subsection{Test Set Performance}

Table~\ref{tab:results} presents the final performance metrics on the
held-out test set, identical for both validation schemes (same training
data and hyperparameters). All three learned models substantially outperform
the dummy baseline ($R^2 = -0.206$). XGBoost achieves the best result on all
three metrics, with Random Forest a close second. The consistency between
walk-forward CV RMSE and test RMSE confirms that the FW scheme provides a
reliable estimate of generalisation performance. Figure~\ref{fig:comparison}
(Appendix) shows side-by-side RMSE bars for both schemes and all models.
Notably, the two schemes produce identical test RMSE values for all
three models (to three decimal places), confirming that the choice
of validation strategy does not affect final model performance when
the test set is held fixed. Because the two schemes yield identical
test-set metrics, all subsequent analyses (feature importance,
predicted vs.\ actual plots) are reported under the TS scheme only;
FW-specific diagnostics are provided in the Appendix.

\begin{table}[htbp]
  \centering
  \small
  \caption{Model performance on the held-out test set (last 20\% by date).
           RPE is on a normalised $[0,1]$ scale. Best value per column
           in bold.}
  \label{tab:results}
  \begin{tabular}{lccc}
    \toprule
    Model & RMSE & MAE & $R^2$ \\
    \midrule
    Dummy baseline & 0.301 & 0.252 & $-$0.206 \\
    Decision Tree  & 0.254 & 0.205 & 0.142 \\
    Random Forest  & 0.242 & 0.193 & 0.217 \\
    XGBoost        & \textbf{0.239} & \textbf{0.189} & \textbf{0.239} \\
    \bottomrule
  \end{tabular}
\end{table}

\subsection{Feature Importance}
\label{subsec:importance}

Figures~\ref{fig:xgb_importance} and~\ref{fig:rf_importance}
display the top-15 features under each importance criterion.
The two models agree on the most informative predictors but diverge
in their ranking of objective load features.

\begin{figure}[H]
  \centering
  \includegraphics[width=0.95\columnwidth]{../figures/fig_xgb_importance_TS.pdf}
  \caption{XGBoost Gain importance (top 15). \emph{Training success
           at $t-2$} dominates by a large margin.}
  \label{fig:xgb_importance}
\end{figure}

Under XGBoost Gain, \emph{perceived training success at $t-2$} (0.195)
and \emph{perceived exertion at $t-2$} (0.173) dominate by a large
margin, with all other features below 0.08. This concentration suggests
that XGBoost's splitting strategy captures most of the predictable
signal through a small number of high-gain splits on subjective
variables. Objective load metrics (total km, km Z5-T1-T2) appear only
below rank 12, with Gain values near 0.01.

\begin{figure}[H]
  \centering
  \includegraphics[width=0.95\columnwidth]{../figures/fig_rf_importance_TS.pdf}
  \caption{Random Forest permutation importance (\%IncMSE, top 15).
           Ranking is more evenly distributed than XGBoost Gain.}
  \label{fig:rf_importance}
\end{figure}

The Random Forest permutation ranking is more distributed:
\emph{total km at $t-1$} (36.9\%) ranks first, followed by
\emph{strength training at $t-1$} (33\%) and
\emph{perceived training success at $t-2$} (28.3\%).
This discrepancy between the two importance measures is expected:
Gain is computed on training splits and can be dominated by a few
high-purity nodes, while permutation importance (\%IncMSE) measures
the effect of each feature on held-out OOB predictions and is
therefore more robust to overfitting artefacts. Both measures agree
that the most recent days ($t-1$, $t-2$) are more informative than
older lags ($t-5$, $t-6$), consistent with the physiological
rationale that acute fatigue decays over two to three days.

The decision tree (Figure~\ref{fig:tree}, Appendix) confirms this:
the root split is on \emph{training success at $t-2$}, and the tree
is shallow (17~leaves), indicating that a small number of conditional
rules on subjective variables capture the main RPE structure.

\subsection{Predicted vs.\ Actual Analysis}
\label{subsec:pred_actual}

Figure~\ref{fig:pred_actual} shows predicted versus actual RPE for
XGBoost on the test set. The cloud aligns broadly with the identity
line, confirming predictive signal. Three patterns emerge:
(i)~under-prediction at high RPE ($> 0.6$), consistent with the
rarity of intense sessions in training data; (ii)~correct
identification of rest days (near-zero load); (iii)~substantial
spread around the diagonal, reflecting irreducible variability in
subjective perception that caps the achievable $R^2$.

\begin{figure}[H]
  \centering
  \includegraphics[width=0.85\columnwidth]{../figures/fig_pred_actual_xgb_TS.pdf}
  \caption{Predicted vs.\ actual RPE for XGBoost on the held-out
           test set. Dashed line: perfect prediction. The spread
           reflects irreducible variability in subjective effort
           perception.}
  \label{fig:pred_actual}
\end{figure}

% -----------------------------------------------------------------------
\section{Discussion}
\label{sec:discussion}

The results demonstrate that tree-based models can predict RPE from recent
training history with moderate accuracy ($R^2 \approx 0.24$ for XGBoost).
The gap between the dummy baseline ($R^2 = -0.206$) and the best model
($R^2 = 0.239$) indicates that the feature set carries genuine predictive
signal, even though a substantial portion of RPE variance remains
unexplained.
This level of accuracy is modest but not unusual: Vallance
et al.~\citep{vallance2023} achieved a 60\% RMSE reduction over a
dummy baseline when predicting RPE in professional soccer players,
though their models included same-session GPS features that are
unavailable in our prediction-ahead setting.

A frank assessment of the $R^2 = 0.239$ achieved by XGBoost is warranted.
The model explains roughly 24\% of RPE variance; over three quarters
remains attributable to factors absent from the dataset, sleep quality,
psychological stress, nutrition, and competitive context. A practitioner
would not modify a training plan on the basis of this model alone.
However, Vallance et al.~\citep{vallance2023} showed that tree-based
models can reduce RPE prediction RMSE by 60\% over a dummy baseline
in professional soccer when same-session GPS data and four weeks of
history are available. Our lower accuracy is expected, since we
deliberately exclude same-session features to enable genuine
one-day-ahead forecasting. The model is therefore better
understood as a complementary screening tool, capable of flagging athletes
at elevated risk of excessive load within a broader multi-dimensional
monitoring system, rather than as a standalone session-level predictor.

The dominance of subjective indicators (perceived training success,
perceived recovery) over objective load metrics in both importance rankings
is consistent with the psychophysiological literature: RPE integrates not
only physiological stress but also motivational and contextual factors that
are partly captured by subjective self-report variables from preceding days
\citep{foster2001}.

Both validation schemes yield the same ranking and near-identical test
metrics, which validates the walk-forward CV as a reliable model-selection
criterion in this temporal setting. The small standard deviation of fold-level
RMSE ($\leq 0.011$) suggests that model performance is stable across
different training windows, i.e., the models do not overfit to any
particular temporal period.

The two-day lag of the strongest predictor (\emph{training success at
$t-2$}) suggests that the psychophysiological response to a suboptimal
session needs more than 24\,h to manifest as elevated perceived
exertion. The weak ranking of objective load metrics reinforces the
case for multi-dimensional monitoring that includes subjective wellness
data alongside GPS-derived external load \citep{foster2001}.

A limitation inherent to decision trees is that predictions in leaves
populated by few training examples carry higher uncertainty, a challenge
addressed by ensemble methods here through bagging (RF) and boosting
(XGBoost) \citep{breiman2001, friedman2001}.

% -----------------------------------------------------------------------
\section{Conclusion}
\label{sec:conclusion}

This study investigated whether the RPE of a competitive runner's
training session can be predicted from six days of training
history using tree-based models. XGBoost achieved the best
performance ($R^2 = 0.239$), followed by Random Forest and
Decision Tree, confirming the expected bias-variance progression
from single tree to bagged and boosted ensembles. Both validation
schemes produced identical rankings, validating walk-forward CV
as a reliable strategy for temporally structured data.

Subjective variables, in particular perceived training success at
$t-2$, consistently outranked objective load metrics, suggesting
that monitoring programmes should prioritise collecting subjective
wellness data alongside GPS-derived metrics.

Several limitations warrant acknowledgement. First, the $R^2$ of
0.239 indicates that a substantial portion of RPE variance remains
unexplained, likely due to factors not captured in the dataset
(sleep quality, stress, nutrition, competitive context). Second,
the prediction horizon is limited to one day ahead; extending to
longer horizons would be of practical interest for weekly training
planning. Third, the model is trained and evaluated on a single
cohort of Dutch competitive runners; generalisation to other
populations (recreational runners, team sport athletes) has not
been tested.

Future work could explore a multi-label formulation that jointly
predicts RPE, perceived recovery, and injury risk, exploiting the
correlations among these outcomes.
Moreover, the subjective variables in the dataset (RPE, training
success, perceived recovery) are inherently imprecise and could be
modelled as uncertain data in the credibilistic sense. The E$^2$M
decision trees proposed by Sutton-Charani~\citep{suttoncharani2014}
would provide a formal framework for integrating this input uncertainty
directly into tree construction, rather than treating it as noise.

% -----------------------------------------------------------------------
% -----------------------------------------------------------------------
\begin{thebibliography}{10}
\small

\bibitem{lovdal2021}
L{\"o}vdal, S., Den Hartigh, R., \& Azzopardi, G. (2021).
Injury prediction in competitive runners with machine learning.
\textit{International Journal of Sports Physiology and Performance},
16(10), 1507--1514.

\bibitem{borg1982}
Borg, G. (1982).
Psychophysical bases of perceived exertion.
\textit{Medicine \& Science in Sports \& Exercise}, 14(5), 377--381.

\bibitem{foster2001}
Foster, C., et al. (2001).
A new approach to monitoring exercise training.
\textit{Journal of Strength and Conditioning Research}, 15(1), 109--115.

\bibitem{vallance2023}
Vallance, E., Sutton-Charani, N., Guyot, P., \& Perrey, S. (2023).
Predictive modeling of the ratings of perceived exertion during training
and competition in professional soccer players.
\textit{Journal of Science and Medicine in Sport}, 26(6), 322--327.

\bibitem{breiman1984}
Breiman, L., Friedman, J., Olshen, R., \& Stone, C. (1984).
\textit{Classification and Regression Trees}. Wadsworth.

\bibitem{suttoncharani2014}
Sutton-Charani, N. (2014).
\textit{Apprentissage \`a partir de donn\'ees et de connaissances
incertaines\,: application \`a la pr\'ediction de la qualit\'e du
caoutchouc} (Doctoral dissertation). Universit\'e de Technologie de
Compi\`egne.

\bibitem{breiman2001}
Breiman, L. (2001).
Random forests.
\textit{Machine Learning}, 45, 5--32.

\bibitem{hastie2009}
Hastie, T., Tibshirani, R., \& Friedman, J. (2009).
\textit{The Elements of Statistical Learning} (2nd ed.). Springer.

\bibitem{friedman2001}
Friedman, J. H. (2001).
Greedy function approximation: a gradient boosting machine.
\textit{Annals of Statistics}, 29(5), 1189--1232.

\bibitem{chen2016}
Chen, T., \& Guestrin, C. (2016).
XGBoost: A scalable tree boosting system.
\textit{Proceedings of KDD 2016}, 785--794.

\end{thebibliography}

\appendix
\onecolumn
\section{Supplementary Figures}
\vspace{-0.3cm}

\begin{figure}[H]
  \centering
  \begin{subfigure}[t]{0.48\textwidth}
    \centering
    \includegraphics[width=\textwidth]{../figures/fig_tree_TS.pdf}
    \caption{TS scheme (17 leaves).}
    \label{fig:tree}
  \end{subfigure}\hfill
  \begin{subfigure}[t]{0.48\textwidth}
    \centering
    \includegraphics[width=\textwidth]{../figures/fig_tree_FW.pdf}
    \caption{FW scheme.}
    \label{fig:tree_fw}
  \end{subfigure}
  \caption{Decision Trees. Root split on \emph{perceived training
           success at $t-2$} in both schemes.}
  \label{fig:trees}
\end{figure}
\vspace{-0.4cm}

\begin{figure}[H]
  \centering
  \begin{subfigure}[t]{0.48\textwidth}
    \centering
    \includegraphics[width=\textwidth]{../figures/fig_xgb_importance_FW.pdf}
    \caption{XGBoost Gain (FW).}
    \label{fig:xgb_importance_fw}
  \end{subfigure}\hfill
  \begin{subfigure}[t]{0.48\textwidth}
    \centering
    \includegraphics[width=\textwidth]{../figures/fig_rf_importance_FW.pdf}
    \caption{RF permutation importance (FW).}
    \label{fig:rf_importance_fw}
  \end{subfigure}
  \caption{Variable importance (top 15) under the FW scheme.
           Rankings mirror the TS results
           (Figures~\ref{fig:xgb_importance}--\ref{fig:rf_importance}).}
  \label{fig:importance_fw}
\end{figure}
\vspace{-0.4cm}

\begin{figure}[H]
  \centering
  \begin{subfigure}[t]{0.48\textwidth}
    \centering
    \includegraphics[width=\textwidth]{../figures/fig_cv_folds_FW.pdf}
    \caption{Walk-forward CV RMSE per fold.}
    \label{fig:cv_folds}
  \end{subfigure}\hfill
  \begin{subfigure}[t]{0.48\textwidth}
    \centering
    \includegraphics[width=\textwidth]{../figures/fig_comparison_TS_vs_FW.pdf}
    \caption{TS vs.\ FW test RMSE comparison.}
    \label{fig:comparison}
  \end{subfigure}
  \caption{Validation scheme analysis.}
  \label{fig:validation}
\end{figure}

\newpage

\begin{figure}[H]
  \centering
  \begin{subfigure}[t]{0.48\textwidth}
    \centering
    \includegraphics[width=\textwidth]{../figures/fig_model_comparison_TS.pdf}
    \caption{TS scheme.}
    \label{fig:model_comparison_ts}
  \end{subfigure}\hfill
  \begin{subfigure}[t]{0.48\textwidth}
    \centering
    \includegraphics[width=\textwidth]{../figures/fig_model_comparison_FW.pdf}
    \caption{FW scheme.}
    \label{fig:model_comparison_fw}
  \end{subfigure}
  \caption{Model comparison (RMSE, MAE, $R^2$) under both validation schemes.}
  \label{fig:model_comparison_appendix}
\end{figure}
\vspace{-0.4cm}

\begin{figure}[H]
  \centering
  \includegraphics[width=0.48\textwidth]{../figures/fig_pred_actual_xgb_FW.pdf}
  \caption{XGBoost predicted vs.\ actual RPE under the FW scheme.}
  \label{fig:pred_actual_fw}
\end{figure}

\end{document}
